% ch12 译文：原文第 6 节 6c（main.tex 第 4216–4655 行）
% [原文 p.45]
\subsection{扭 Künneth 公式}

本节以 Künneth 定理收尾。设 $X$ 与 $Y$ 分别是 $G$-概形与 $H$-概形，且分别存在由 $G$-不变、$H$-不变的仿射开子概形组成的开覆盖 $\mathfrak{U}$、$\mathfrak{V}$。回忆：$X\times Y$ 有由形如 $U\times V\in\mathfrak{U}\times\mathfrak{V}$ 的仿射概形组成的开覆盖，这显然是 $G\times H$-不变的。于是我们可以定义等变 Künneth 函子
$$
\Perf_G(X;\mathfrak{U})_{\dg}
\otimes
\Perf_H(Y;\mathfrak{V})_{\dg}
\xrightarrow{\Ku}
\Perf_{G\times H}(X\times Y;\mathfrak{U}\times\mathfrak{V})_{\dg},
$$
它在 dg 加强上的定义如下：
\begin{defn}
$\Ku$ 是一个函子，其定义如下：
\begin{itemize}
\item
对象：
它把对象
$
\big((\mathcal{F}^{\bullet},\theta^{\bullet}),(\mathcal{G}^{\bullet},\eta^{\bullet})\big)
$
送到对象
$$
\Big(
\bigoplus_{n\in\mathbb{Z}}
\big(
(\mathcal{F}^{\bullet}\boxtimes\mathcal{G}^{\bullet})^n,
(\theta^{\bullet}\boxtimes\eta^{\bullet})^n
\big),
\
d_{\mathcal{F}^{\bullet}\boxtimes\mathcal{G}^{\bullet}}
:=
\big(
d_{\mathcal{F}^{\bullet}\boxtimes\mathcal{G}^{\bullet}}^n
:
(\mathcal{F}^{\bullet}\boxtimes\mathcal{G}^{\bullet})^n
\rightarrow
(\mathcal{F}^{\bullet}\boxtimes\mathcal{G}^{\bullet})^{n+1}
\big)_{n\in\mathbb{Z}}
\Big)
$$
（见定义 \ref{dfn:GEquivQCohdg}），其中
$$
\big(
(\mathcal{F}^{\bullet}\boxtimes\mathcal{G}^{\bullet})^n,
(\theta^{\bullet}\boxtimes\eta^{\bullet})^n
\big)
:=
\bigoplus_{i+j=n}
(
\mathcal{F}^i\boxtimes\mathcal{G}^j,
\
\theta^i\boxtimes\eta^j
)
\quad\text{对每个 }n\in\mathbb{Z},
$$
（注意 $
(\mathcal{F}^i,\theta^i)
\boxtimes
(\mathcal{G}^j,\eta^j)
=
(
\mathcal{F}^i\boxtimes\mathcal{G}^j,
\
\theta^i\boxtimes\eta^j
)
\in\Mod_{G\times H}(X\times Y)
$
的定义见定义 \ref{defn:ExternalTensorProductOfEquivariantModules}），而微分在每个直和项 $\mathcal{F}^i\boxtimes\mathcal{G}^j$（其中 $i+j=n$）上由下式确定：
\[
d_{\mathcal F^\bullet\boxtimes\mathcal G^\bullet}^n
\big|_{\mathcal F^i\boxtimes\mathcal G^j}
:=
d_{\mathcal F}^i\boxtimes\id_{\mathcal G^j}
+
(-1)^i\id_{\mathcal F^i}\boxtimes d_{\mathcal G}^j
:
\mathcal{F}^i\boxtimes\mathcal{G}^j
\rightarrow
(\mathcal{F}^{i+1}\boxtimes\mathcal{G}^j)
\oplus
(\mathcal{F}^i\boxtimes\mathcal{G}^{j+1}).
\]

% [原文 p.46]
\item
态射：
对齐次态射 $f\in C_{\otimes}^{p+q}$ 与 $g\in C_{\otimes}^{r+s}$，$f\otimes g$ 的像为 $T^{\ast}(f\otimes g)\in C_{\times}^{p+q+r+s}$，其中 $T^{\ast}$ 是第 \ref{subsec:Eilenberg-Zilber-contraction-Cech} 节定义的映射。
\item
高阶结构：我们定义第 $n$ 个 Künneth 函子
\begin{align*}
\Ku^n\colon &\hom_{\Perf (X;\mathfrak{U})_{dg}\otimes\Perf(Y;\mathfrak{V})_{dg}}\left(\left(
    (\mathcal{F}_{n-1}^{\bullet},\theta_{n-1}^{\bullet}),(\mathcal{G}_{n-1}^{\bullet},\eta_{n-1}^{\bullet})\right),
    \left((\mathcal{F}_n^{\bullet},\theta_n^{\bullet}),(\mathcal{G}_{n}^{\bullet},\eta_n^{\bullet})\right)\right)
    \otimes\dotsb\\ 
    \otimes
    &\hom_{\Perf (X;\mathfrak{U})_{dg}\otimes\Perf(Y;\mathfrak{V})_{dg}}\left(\left(
    (\mathcal{F}_0^{\bullet},\theta_0^{\bullet}),(\mathcal{G}_0^{\bullet},\eta_0^{\bullet})\right),
    \left((\mathcal{F}_1^{\bullet},\theta_1^{\bullet}),(\mathcal{G}_1^{\bullet},\eta_1^{\bullet})\right)\right)\\
    \to &\hom_{\Perf_{G\times H} (X\times Y;\mathfrak{U}\times\mathfrak{V})_{dg}}
    \left((\mathcal{F}_0^{\bullet},\theta_0^{\bullet})\boxtimes (\mathcal{G}_0^{\bullet},\eta_0^{\bullet}),(\mathcal{F}_n^{\bullet},\theta_n^{\bullet})\boxtimes (\mathcal{G}_n^{\bullet},\eta_n^{\bullet})\right)
\end{align*}
其公式为
\[
\resizebox{\linewidth}{!}{$\displaystyle
    \Ku^n (f_{n-1}\otimes g_{n-1},\dotsb, f_0\otimes g_0)=\sum_{\ell=1}^{n-1} (-1)^{(n-\ell -1)(t_{\ell}'+u_{\ell}'-1)}\mathbb{H}^{\ast}\left(
        \Ku^{n-\ell} (f_{n-1}\otimes g_{n-1},\dotsb, f_{\ell}\otimes g_{\ell})\circ\Ku^{\ell} (f_0\otimes g_0,\dotsb, f_{\ell-1}\otimes g_{\ell-1})
        \right),$}
\]
其中算子 $\mathbb{H}^{\ast}$ 的定义见附录 \ref{sec:Eilenberg-Zilber-constraction}。
\end{itemize}
\end{defn}

我们可以直接验证：
\begin{prp}
族 $(\Ku^n)_{n\in\mathbb{Z}_{\geq0}}$ 构成从 $\Perf_G (X;\mathfrak{U})_{dg}\otimes\Perf_H (Y;\mathfrak{V})_{dg}$ 到 $\Perf (X\times Y;\mathfrak{U}\times\mathfrak{V})_{dg}$ 的 $A_{\infty}$-函子。
\end{prp}
\begin{proof}
完整的验证相当冗长，我们只给出第一个非平凡关系的验证，以展示其中的模式。由附录 \ref{sec:Eilenberg-Zilber-constraction} 可知 $T^{\ast} S^{\ast} =\id$ 且 $S^{\ast} T^{\ast} -\id =\mathbb{H}^{\ast}d+d\mathbb{H}^{\ast}$，从而
\begin{align*}
d\Ku^2& (f_1\otimes g_1,f_0\otimes g_0)=d\mathbb{H}^{\ast} (\Ku (f_1\otimes g_1)\circ\Ku (f_0\otimes g_0))\\ 
&=-\mathbb{H}^{\ast} (\Ku (d(f_1\otimes g_1))\circ \Ku (f_0\otimes g_0)) -(-1)^{\vert f_0\otimes g_0\vert }\mathbb{H}^{\ast} (\Ku(f_1\otimes g_1)\circ\Ku (d(f_0\otimes g_0)))\\
&+T^{\ast} (S^{\ast}(\Ku (f_1\otimes g_1))\circ S^{\ast} (\Ku (f_0\otimes g_0))) -\Ku (f_1\otimes g_1)\circ\Ku (f_0\otimes g_0)\\ 
&=-(-1)^{\vert f_0\otimes g_0\vert }\Ku^2 (d(f_1\otimes g_1),f_0\otimes g_0) - \Ku^2(f_1\otimes g_1,d(f_0\otimes g_0))\\ 
&+\Ku((f_1\otimes g_1)\circ (f_0\otimes g_0))-\Ku (f_1\otimes g_1)\circ\Ku (f_0\otimes g_0),
\end{align*}
而这正是 $\Ku^2$ 的 $A_{\infty}$-关系。
\end{proof}

现在可以证明 Künneth 嵌入。
\begin{thm}\label{thm:Kunneth-embedding}
预先给定的函子
$$
\Ku\colon\Perf_G(X;\mathfrak{U})_{\dg}
\otimes
\Perf_H(Y;\mathfrak{V})_{\dg}
\hookrightarrow
\Perf_{G\times H}(X\times Y;\mathfrak{U}\times\mathfrak{V})_{\dg},
$$
$$
(\mathcal{F},\mathcal{G})
\quad\mapsto\quad
\mathcal{F}\boxtimes\mathcal{G}
~:=~
\pi_X^*\mathcal{F}
\underset{\mathcal{O}_{X\times Y}}{\otimes}
\pi_Y^*\mathcal{G},
$$
$$
f\otimes g
\quad\mapsto\quad
f\boxtimes g
~:=~
\pi_X^*f
\otimes
\pi_Y^*g,
$$
是 dg-范畴的拟嵌入。
\end{thm}
% [原文 p.47]
\begin{proof}
由附录 \ref{sec:Eilenberg-Zilber-constraction} 我们知道 $\Ku^1 =T^{\ast}$ 是链复形的拟等价，其拟逆为 $S^{\ast}$；因此过渡到同调范畴后，$H^0 (\Ku)\colon H^0 (\Perf_G (X;\mathfrak{U})_{dg}\otimes\Perf_H (Y;\mathfrak{V})_{dg})\to H^0 (\Perf_{G\times H} (X\times Y;\mathfrak{U}\times\mathfrak{V})_{dg})$ 是范畴的嵌入。这就证明了 $\Ku$ 是拟嵌入。
\end{proof}

\begin{crl}
$\Ku$ 诱导拟嵌入
$$
\Perf_G(X)_{\dg}
\underset{\widehat{G}}{\otimes}
\Perf_G(Y)_{\dg}
\hookrightarrow
\Perf_G(X\times Y)_{\dg},
$$
$$
(\mathcal{F},\mathcal{G})
\quad\mapsto\quad
\mathcal{F}\boxtimes\mathcal{G}
~:=~
\pi_X^*\mathcal{F}
\underset{\mathcal{O}_{X\times Y}}{\otimes}
\pi_Y^*\mathcal{G},
$$
$$
f\otimes g
\quad\mapsto\quad
f\boxtimes g
~:=~
\pi_X^*f
\otimes
\pi_Y^*g.
$$
\end{crl}
\begin{proof}
由定理 \ref{thm:Kunneth-embedding} 得到拟嵌入 $\Perf_G (X)_{dg}\otimes\Perf_G(Y)_{dg}\hookrightarrow\Perf_{G\times G} (X\times Y)_{dg}$。另一方面，考虑对角嵌入 $\Delta\colon G\hookrightarrow G\times G$，其正交子群作用 $\widehat{G}\cong (\Delta (G))^{\perp}\curvearrowright\Perf_{G\times G} (Y)$ 保持子范畴 $\Perf_G (X)_{dg}\otimes\Perf_G(Y)_{dg}$，因此由推论 \ref{crl:RestrictionFunctionEm} 与命题 \ref{prp:twisted-tensor-equal-quotient} 得到拟嵌入
\begin{align*}
    \Perf_G (X)_{dg}\underset{\widehat{G}}{\otimes}\Perf_G(Y)_{dg}&\simeq\left(\Perf_G(X)\otimes\Perf_G(Y)_{dg}\right)\big/\widehat{G}\\ 
    &\hookrightarrow\Perf_{G\times G} (X\times Y)_{dg}\big/\widehat{G}\hookrightarrow\Perf_G (X\times Y)_{dg}.\qedhere
\end{align*}
\end{proof}

\begin{yizhu}[原文此处的断句与记号小疵]
原文在「$X\times Y$ 的仿射开覆盖」一句处以逗号接「Then we can define」，句子并不完整；译文按「这显然是 $G\times H$-不变的」收句，再另起一句。此外，定义中第 $n$ 个 Künneth 函子的定义域写成 $\Perf(X;\mathfrak{U})_{dg}\otimes\Perf(Y;\mathfrak{V})_{dg}$，漏掉了下标 $G$、$H$，而值域却写作 $\Perf_{G\times H}(\cdots)_{dg}$；译文照录原文，读者按上下文补足即可。推论证明中「$\widehat{G}\cong(\Delta(G))^\perp$ 作用在 $\Perf_{G\times G}(Y)$ 上」也应是 $\Perf_{G\times G}(X\times Y)$。
\end{yizhu}

\begin{yizhu}[「扭」扭在哪里]
先忘掉群作用，Künneth 嵌入说的是 $\Perf(X)\otimes\Perf(Y)\hookrightarrow\Perf(X\times Y)$。但在等变情形，$\Perf_G(X)\otimes\Perf_G(Y)$ 只给出 $\Perf_{G\times G}(X\times Y)$，并不是 $\Perf_G(X\times Y)$。补救办法是沿对角嵌入 $\Delta\colon G\hookrightarrow G\times G$ 把 $G\times G$-作用限制为 $G$-作用，再商掉正交子群 $(\Delta(G))^\perp\cong\widehat{G}$ 的作用——这一步正是「扭」：张量积要换成带 $\widehat{G}$-修正的版本 $\underset{\widehat{G}}{\otimes}$，只保留与 $\widehat{G}$-作用相容的那部分。推论 \ref{crl:RestrictionFunctionEm} 提供的轨道范畴嵌入，就是完成这次商的技术保证。
\end{yizhu}

\begin{yizhu}[为什么给出的是 $A_{\infty}$-函子而不是 dg-函子]
收缩只保证 $T^{\ast}$ 是链映射、$S^{\ast}T^{\ast}-\id=\mathbb{H}^{\ast}d+d\mathbb{H}^{\ast}$ 只在同伦意义下成立，所以 $T^{\ast}$ 一般并不严格保持乘法：$\Ku(f\circ f')$ 与 $\Ku(f)\circ\Ku(f')$ 只差一个同伦。把这些同伦修正逐阶打包，就得到高阶分量 $\Ku^n$，它们整体构成 $A_{\infty}$-函子；上文的验证只需检验第一个非平凡关系 $\Ku^2$，其余关系由同一模式归纳得到。
\end{yizhu}

\begin{conceptbox}{扭 Künneth 公式}
Künneth 公式说，两个空间（上）同调的张量积给出乘积空间的（上）同调；在范畴层面，它表现为 $\Perf(X)\otimes\Perf(Y)\simeq\Perf(X\times Y)$ 型的等价。当两个空间各自带群作用、而我们只想在等变范畴里做张量积时，普通的张量积会多出一个「多余」的群方向，必须沿对角嵌入把作用拉回并商掉正交子群 $\widehat{G}$ 的作用，这就是扭 Künneth 公式：张量积换成 $\underset{\widehat{G}}{\otimes}$，结论仍是一个拟嵌入。本文用 dg-范畴版本的 Eilenberg–Zilber 收缩来构造并验证这个嵌入，它是第 6 节把 $G$-等变层与 Fukaya 范畴的 Künneth 对接起来的技术核心。
\par\smallskip\noindent{\footnotesize 参考：nLab，\emph{Künneth theorem}（ncatlab.org/nlab/show/Künneth+theorem）；Eilenberg–Zilber, \emph{On products of complexes}, Amer. J. Math. 75 (1953)；本文第 6 节 6c 与附录 A。}
\end{conceptbox}

\begin{keytakeaways}
\item 第 6 节把凝聚层的导出范畴提升为 dg 范畴：用 $\mathfrak{U}$ 上的 Čech 复形定义 $\Coh(X;\mathfrak{U})_{\dg}$，等变情形则把 Čech 复形换成 $G$-不变版本，得到 $\Coh_G(X;\mathfrak{U})_{\dg}$。
\item 对闭子群概形 $H\subseteq G$，限制 dg-函子 $\Res_G^H$ 诱导轨道范畴的拟嵌入 $\Coh_G(X;\mathfrak{U})_{\dg}/H^\perp\hookrightarrow\Coh_H(X;\mathfrak{U})_{\dg}$；这是「扭」在代数侧的来源。
\item 扭 Künneth 公式断言 $\Perf_G(X)_{\dg}\underset{\widehat{G}}{\otimes}\Perf_G(Y)_{\dg}\hookrightarrow\Perf_G(X\times Y)_{\dg}$ 是拟嵌入，其构造依赖 Eilenberg–Zilber 收缩给出的同伦逆与高阶修正。
\end{keytakeaways}
